\thispagestyle{plain}

% settings for abstract
\setstretch{1.3}
\setlength{\parskip}{12pt}

% title
{\fontsize{14}{18}\selectfont\textbf{چکیده}}

\vspace{24pt} % distance

% abstract 
{\fontsize{13}{16}\selectfont
چکیده بخشی از پایان‌نامه است که خواننده را به مطالعه آن علاقمند می‌کند و یا از آن می‌گریزاند.
چکیده باید ترجیحاً در یک صفحه باشد (تقریباً تمامی چکیده پایان‌نامه‌ها در یک صفحه قابل نگارش است).
در نگارش چکیده نکات زیر باید رعایت شود. متن چکیده باید مزین به کلمه‌ها و عبارات سلیس، آشنا، با
معنی و روشن باشد. بگونه ای که با حدود 300 تا 500 کلمه بتواند خواننده را به خواندن پایان‌نامه راغب
نماید. چکیده، جدا از پایان‌نامه باید به تنهایی گویا و مستقل باشد. در چکیده باید از ذکر منابع، اشاره به
جداول و نمودارها اجتناب شود. تمیز بودن مطلب، نداشتن غلطهای املایی یا دستور زبانی و رعایت دقت و
تسلسل روند نگارش چکیده از نکات مهم دیگری است که باید درنظر گرفته شود. در چکیده پایان‌نامه باید از
درج مشخصات مربوط به پایان‌نامه خودداری شود. کلمات کلیدی در انتهای چکیده فارسی و انگلیسی آورده
شود. محتوای چکیده‌ها بر اساس موضوع و گرایش تحقیق طبقه‌بندی می‌شود و به همین جهت وجود کلمات
شاخص و کلیدی، مراکز اطلاعاتی را در طبقه‌بندی دقیق و سریع پایان‌نامه یاری می‌دهد. کلمات کلیدی،
راهنمای نکات مهم موجود در پایان‌نامه هستند. بنابراین باید در حد امکان کلمه‌ها یا عباراتی انتخاب شد که
ماهیت، محتوا و گرایش کار را به وضوح روشن نماید. چکیده باید منعکس کننده اصل موضوع باشد. در
چکیده باید اهداف تحقیق مورد توجه قرار گیرد. تأکید روی اطلاعات تازه (یافته‌ها) و اصطلاحات جدید یا
نظریه‌ها، فرضیه‌ها، نتایج و پیشنهادها متمرکز شود. اگر در پایان‌نامه روش نوینی برای اولین بار ارائه می‌شود
و تا به حال معمول نبوده است، با جزئیات بیشتری ذکر شود. شایان ذکر است چکیده فارسی و انگلیسی
باید حتماً به تأیید استاد راهنما رسیده باشد.
}

\vspace{12pt}

% keywords
{\fontsize{13}{16}\selectfont\textbf{واژه‌های کلیدی:}}
{\fontsize{13}{16}\selectfont
تعداد کلمات یا عبارات کلیدی حداکثر می‌تواند پنج کلمه یا عبارت باشد. کلید واژه‌ها
باید با واژه‌های اصلی عنوان و مسئله تحقیق تناسب داشته باشند.
}
